% GENERATED FILE. Edit canonical lessons, then run npm run generate:books.
% canonical-source-hash: fnv1a64:bbfb14223b5b87b5
% canonical-lessons: LA-C500-cupido, LA-C500-cupidus, LA-C500-avidus, LA-C500-invideo, LA-C500-rivalis

\chapter{Desire and envy}
\label{ch:la-desire-and-envy}
\hlchaptermodality{500}

\begin{chapteropening}
\textbf{By the end of this chapter:} \emph{I can talk about desire, greed and envy.}

Complete the last lesson of chapter 500: I can talk about desire, greed and envy.
\end{chapteropening}

\section[cupīdō, cupīdinis]{cupīdō, cupīdinis — desire, longing}

\label{lesson:LA-C500-cupido}



\begin{quote}
Before the new one: say the Latin for to amuse, then the Latin for to please.
\end{quote}

\subsection*{You'll want to know: cupīdō, cupīdinis}
\textbf{cupīdō, cupīdinis} — \textquotedblleft{}desire, longing\textquotedblright{}.

\begin{quote}
Cupīdō pecūniae multōs hominēs malōs facit.
\end{quote}

\begin{quote}
\emph{Desire for money makes many people bad.}
\end{quote}

Feminine in prose, third declension; as the god of love it is masculine. Built on \emph{cupiō}.

\begin{cousinweb}
\textbf{cupīdō} is built on \emph{cupiō} \textquotedblleft{}to desire\textquotedblright{} (EtymDB; Lewis \& Short); \emph{cupiditās}, already known, comes from the same verb, and English \emph{cupidity} from that.
\end{cousinweb}

\begin{grammarlens}[title={What you've built}]
One more word for telling, explaining and describing at B1.
\end{grammarlens}

\subsection*{Your turn}
\begin{itemize}
\raggedright
  \item \emph{Say it:} \emph{cupīdō, cupīdinis}
  \item \emph{Say it:} \emph{cupīdō, cupīdinis}, once more
  \item \emph{Say it:} a sentence of your own with \emph{cupīdō, cupīdinis}
  \item \emph{Recall:} say the Latin for to amuse, then the Latin for to please, then say \emph{cupīdō, cupīdinis} again
\end{itemize}

\begin{tcolorbox}[breakable,colback=teal!4,colframe=teal!35!black,title={Before you move on}]
What does cupīdō, cupīdinis mean? (\textquotedblleft{}Desire, longing\textquotedblright{}.) Say it once more.
\end{tcolorbox}
\section[cupidus, cupida, cupidum]{cupidus, cupida, cupidum — eager, desirous}

\label{lesson:LA-C500-cupidus}



\begin{quote}
Before the new one: say the Latin for to please, then the Latin for desire.
\end{quote}

\subsection*{You'll want to know: cupidus, cupida, cupidum}
\textbf{cupidus, cupida, cupidum} — \textquotedblleft{}eager, desirous\textquotedblright{}.

\begin{quote}
Discipulus cupidus litterārum est.
\end{quote}

\begin{quote}
\emph{The pupil is eager for learning.}
\end{quote}

With the genitive: \emph{cupidus pecūniae} — greedy for money.

\begin{cousinweb}
\textbf{cupidus} is built on \emph{cupiō} \textquotedblleft{}to desire\textquotedblright{} with the adjective ending \emph{-idus} (Lewis \& Short; EtymDB).
\end{cousinweb}

\begin{grammarlens}[title={What you've built}]
One more word for telling, explaining and describing at B1.
\end{grammarlens}

\subsection*{Your turn}
\begin{itemize}
\raggedright
  \item \emph{Say it:} \emph{cupidus, cupida, cupidum}
  \item \emph{Say it:} \emph{cupidus, cupida, cupidum}, once more
  \item \emph{Say it:} a sentence of your own with \emph{cupidus, cupida, cupidum}
  \item \emph{Recall:} say the Latin for to please, then the Latin for desire, then say \emph{cupidus, cupida, cupidum} again
\end{itemize}

\begin{tcolorbox}[breakable,colback=teal!4,colframe=teal!35!black,title={Before you move on}]
What does cupidus, cupida, cupidum mean? (\textquotedblleft{}Eager, desirous\textquotedblright{}.) Say it once more.
\end{tcolorbox}
\section[avidus, avida, avidum]{avidus, avida, avidum — greedy, eager}

\label{lesson:LA-C500-avidus}



\begin{quote}
Before the new one: say the Latin for desire, then the Latin for eager.
\end{quote}

\subsection*{You'll want to know: avidus, avida, avidum}
\textbf{avidus, avida, avidum} — \textquotedblleft{}greedy, eager\textquotedblright{}.

\begin{quote}
Canis avidus cibum statim edit.
\end{quote}

\begin{quote}
\emph{The greedy dog eats the food at once.}
\end{quote}

With the genitive: \emph{avidus glōriae} — hungry for glory. English \emph{avid}.

\begin{cousinweb}
\textbf{avidus} is built on \emph{aveō} \textquotedblleft{}to long for\textquotedblright{} (Lewis \& Short; EtymDB). English \emph{avid} comes from it.
\end{cousinweb}

\begin{grammarlens}[title={What you've built}]
One more word for telling, explaining and describing at B1.
\end{grammarlens}

\subsection*{Your turn}
\begin{itemize}
\raggedright
  \item \emph{Say it:} \emph{avidus, avida, avidum}
  \item \emph{Say it:} \emph{avidus, avida, avidum}, once more
  \item \emph{Say it:} a sentence of your own with \emph{avidus, avida, avidum}
  \item \emph{Recall:} say the Latin for desire, then the Latin for eager, then say \emph{avidus, avida, avidum} again
\end{itemize}

\begin{tcolorbox}[breakable,colback=teal!4,colframe=teal!35!black,title={Before you move on}]
What does avidus, avida, avidum mean? (\textquotedblleft{}Greedy, eager\textquotedblright{}.) Say it once more.
\end{tcolorbox}
\section[invideō, invidēre]{invideō, invidēre — to envy, begrudge}

\label{lesson:LA-C500-invideo}



\begin{quote}
Before the new one: say the Latin for eager, then the Latin for greedy.
\end{quote}

\subsection*{You'll want to know: invideō, invidēre}
\textbf{invideō, invidēre} — \textquotedblleft{}to envy, begrudge\textquotedblright{}.

\begin{quote}
Soror sorōrī novam stolam invidet.
\end{quote}

\begin{quote}
\emph{One sister envies the other her new dress.}
\end{quote}

\emph{invideō, invidēre, invīdī, invīsum}, with the dative of the person: \emph{tibi invideō} — I envy you. \emph{Invidia} and \emph{invidus}, already known, come from it.

\begin{cousinweb}
\textbf{invideō} is built on \emph{in-} \textquotedblleft{}upon, against\textquotedblright{} and \emph{videō} \textquotedblleft{}to see\textquotedblright{} (Lewis \& Short), first \textquotedblleft{}to cast an evil eye on\textquotedblright{}. English \emph{envy} comes from \emph{invidia} through Old French.
\end{cousinweb}

\begin{grammarlens}[title={What you've built}]
One more word for telling, explaining and describing at B1.
\end{grammarlens}

\subsection*{Your turn}
\begin{itemize}
\raggedright
  \item \emph{Say it:} \emph{invideō, invidēre}
  \item \emph{Say it:} \emph{invideō, invidēre}, once more
  \item \emph{Say it:} a sentence of your own with \emph{invideō, invidēre}
  \item \emph{Recall:} say the Latin for eager, then the Latin for greedy, then say \emph{invideō, invidēre} again
\end{itemize}

\begin{tcolorbox}[breakable,colback=teal!4,colframe=teal!35!black,title={Before you move on}]
What does invideō, invidēre mean? (\textquotedblleft{}To envy, begrudge\textquotedblright{}.) Say it once more.
\end{tcolorbox}
\section[rīvālis, rīvālis]{rīvālis, rīvālis — a rival (in love)}

\label{lesson:LA-C500-rivalis}



\begin{quote}
Before the new one: say the Latin for greedy, then the Latin for to envy.
\end{quote}

\subsection*{You'll want to know: rīvālis, rīvālis}
\textbf{rīvālis, rīvālis} — \textquotedblleft{}a rival (in love)\textquotedblright{}.

\begin{quote}
Duo iuvenēs eandem puellam amant; rīvālēs sunt.
\end{quote}

\begin{quote}
\emph{Two young men love the same girl; they are rivals.}
\end{quote}

Masculine, third declension. First someone who shares a stream with you, then a rival in love. English \emph{rival}.

\begin{cousinweb}
\textbf{rīvālis} is built on \emph{rīvus} \textquotedblleft{}a stream\textquotedblright{} (Lewis \& Short; EtymDB): first those who used the same brook, then rivals in love. English \emph{rival} comes from it.
\end{cousinweb}

\begin{grammarlens}[title={What you've built}]
One more word for telling, explaining and describing at B1.
\end{grammarlens}

\subsection*{Your turn}
\begin{itemize}
\raggedright
  \item \emph{Say it:} \emph{rīvālis, rīvālis}
  \item \emph{Say it:} \emph{rīvālis, rīvālis}, once more
  \item \emph{Say it:} a sentence of your own with \emph{rīvālis, rīvālis}
  \item \emph{Recall:} say the Latin for greedy, then the Latin for to envy, then say \emph{rīvālis, rīvālis} again
\end{itemize}

\begin{tcolorbox}[breakable,colback=teal!4,colframe=teal!35!black,title={Before you move on}]
What does rīvālis, rīvālis mean? (\textquotedblleft{}A rival (in love)\textquotedblright{}.) Say it once more.
\end{tcolorbox}
